\documentclass{report}

\usepackage{graphicx}
\usepackage{amsmath}
\usepackage{booktabs}
\usepackage[backend=biber, style=numeric]{biblatex}


\begin{document}

\subsection{GMOS Readout Circuit}

\subsubsection{Measurement Principle}

Methods for biasing and signal readout of CMOS-based thermal sensors were
explored in detail in~\cite{TODO}. That work distinguishes two readout
topologies: a diode-connected (2T) configuration and an amplification (3T)
configuration, and derives the fundamental equations governing the electrical
characteristics of the GMOS device. In our setup we adopt the 2T connection
for its simplicity and reproducibility.

In the 2T configuration the GMOS drain is shorted to its gate, forcing
$V_{DS} = V_{GS}$. For $V_{GS} > V_{th}$ the transistor is in saturation,
and in the subthreshold regime the drain current is

\begin{equation}
  I_D = I_0 \exp\!\left(\frac{V_{GS} - V_{th}}{n V_T}\right),
  \label{eq:subthreshold}
\end{equation}

where $I_0$ is a process-dependent pre-exponential current, $V_{th}$ is the
threshold voltage, $n$ is the subthreshold ideality factor, and
$V_T = k_B T / q$ is the thermal voltage. Because the GMOS threshold voltage
shifts with gas concentration via the catalytic reaction at the gate, the
drain current encodes the chemical signal.

\subsubsection{Differential Readout and Common-Mode Rejection}

In analogy with the TMOS bolometer, the GMOS performs thermal-offset
subtraction by reading the differential drain voltage between a blind and an
active pixel. The active pixel is exposed to the analyte gas; the blind pixel
is identical but sealed or passivated. Both pixels experience the same
ambient temperature, supply noise, and heater-induced thermal background.
Taking their difference therefore isolates only the additional
threshold-voltage shift caused by the gas reaction:

\begin{equation}
  \Delta V_{\mathrm{out}} = V_A - V_B
    = \bigl(V_{DD} - I_{D,\mathrm{act}} R_s\bigr)
    - \bigl(V_{DD} - I_{D,\mathrm{bld}} R_s\bigr)
    = \bigl(I_{D,\mathrm{bld}} - I_{D,\mathrm{act}}\bigr) R_s,
  \label{eq:differential}
\end{equation}

where $V_A$ and $V_B$ are the drain voltages of the active and blind pixels,
$R_s$ is the series bias resistor, and $V_{DD}$ is the supply voltage from
the B2961A. The LIA's differential input mode (A$-$B) performs this
subtraction in hardware, further suppressing any residual common-mode noise.

\subsubsection{Bias Resistor and Current-Source Approximation}

A large series resistor $R_s$ placed between the B2961A supply and each drain
approximates a constant current source. For this approximation to hold the
resistor must dominate over the transistor's small-signal output resistance:

\begin{equation}
  R_s \gg r_{ds} = \left.\frac{\partial V_{DS}}{\partial I_D}\right|_{Q}.
  \label{eq:current_source_cond}
\end{equation}

Additionally, when a chemical reaction shifts $V_{th}$ and hence $I_D$, the
resulting current variation must be small enough that the operating point does
not move significantly. Using Eq.~\eqref{eq:subthreshold}, a threshold shift
of $\delta V_{th}$ produces a fractional current change
$\delta I_D / I_D = \delta V_{th} / (n V_T)$.
Requiring this to remain small imposes

\begin{equation}
  R_s \gg \frac{n V_T}{I_D},
  \label{eq:bias_stability}
\end{equation}

which is satisfied by using $R_s \geq 500\,\Omega$ at typical GMOS
subthreshold bias levels. Both blind and active bias resistors must be of
equal value so that Eq.~\eqref{eq:differential} is not corrupted by a
resistor mismatch.

\subsubsection{AC Heater Excitation and Lock-In Detection}

\paragraph{Heater drive.}
Rather than applying a DC bias to the integrated heater resistors, the LIA
reference output applies a sinusoidal voltage at frequency $f_{\mathrm{ref}}$:

\begin{equation}
  V_H(t) = V_{\mathrm{ref}} \sin(2\pi f_{\mathrm{ref}}\, t).
  \label{eq:heater_drive}
\end{equation}

The instantaneous power dissipated in each heater is

\begin{equation}
  P_H(t) = \frac{V_H^2(t)}{R_H}
          = \frac{V_{\mathrm{ref}}^2}{2 R_H}
            \Bigl[1 - \cos(4\pi f_{\mathrm{ref}}\, t)\Bigr],
  \label{eq:heater_power}
\end{equation}

so the thermal modulation of the transistor occurs at twice the drive
frequency $2f_{\mathrm{ref}}$.
% TODO: confirm whether the LIA reference is locked to f_ref or 2*f_ref.
Both heater resistors, $R_{H,\mathrm{act}}$ and $R_{H,\mathrm{bld}}$, are
driven in parallel and are assumed to have equal values and equal temperature
coefficients of resistance (TCR). Imbalances between the two heater resistors
will be addressed separately.
% TODO: add reference to heater imbalance correction section.

\paragraph{Lock-in detection.}
The LIA demodulates the differential drain signal $\Delta V_{\mathrm{out}}(t)$
at the reference frequency, extracting the in-phase ($X$) and quadrature ($Y$)
components:

\begin{align}
  X &= \frac{2}{T}\int_0^T \Delta V_{\mathrm{out}}(t)\,
        \sin(2\pi f_{\mathrm{ref}}\, t)\, \mathrm{d}t,
  \label{eq:lia_x} \\
  Y &= \frac{2}{T}\int_0^T \Delta V_{\mathrm{out}}(t)\,
        \cos(2\pi f_{\mathrm{ref}}\, t)\, \mathrm{d}t.
  \label{eq:lia_y}
\end{align}

The magnitude $R = \sqrt{X^2 + Y^2}$ is the primary readout quantity and
is correlated with gas concentration. Using AC excitation and narrow-band
demodulation moves the signal of interest away from DC, where flicker
($1/f$) noise is largest, to a frequency band where the amplifier noise
floor is much lower.

\subsubsection{Noise Considerations}

The dominant extrinsic noise source in this circuit is the Johnson (thermal)
noise of the bias resistor $R_s$. Its one-sided voltage spectral density is

\begin{equation}
  S_{v,R_s}(f) = 4 k_B T R_s,
  \label{eq:johnson}
\end{equation}

giving an RMS noise voltage over the LIA's effective noise bandwidth
$\Delta f_{\mathrm{LIA}}$:

\begin{equation}
  v_{n,R_s} = \sqrt{4 k_B T R_s\, \Delta f_{\mathrm{LIA}}}.
  \label{eq:johnson_rms}
\end{equation}

The LIA's effective noise bandwidth is determined by its time constant $\tau$
and filter order; for a first-order filter
$\Delta f_{\mathrm{LIA}} = 1/(4\tau)$, which can be reduced to sub-hertz
values by increasing $\tau$. This narrow bandwidth is the primary mechanism
by which the lock-in technique achieves a high signal-to-noise ratio.

% TODO: add shot noise of MOSFET drain current if relevant at these bias
%       levels.

\subsection{Circuit Configuration and Instruments}

\subsubsection{Required Instruments}

\begin{itemize}
  \item SR860 Lock-In Amplifier.
  \item HP 8116A Signal Generator.
  \item HP 6624A DC Power Supply.
  \item Keysight B2961A Low Noise Power Supply.
\end{itemize}

\subsubsection{Electrical Schematic}

Figure~\ref{fig:gmos_2t_schematic} shows the 2T differential readout
connection for the active and blind GMOS pair. Each NMOS has its drain
shorted to its gate (2T configuration). A series bias resistor $R_s$
connects each drain to the B2961A supply $V_{DD}$, and the drain nodes feed
the LIA differential input channels A and B. The heater resistors
$R_{H,\mathrm{act}}$ and $R_{H,\mathrm{bld}}$ are connected in parallel and
driven by the LIA sine reference output.

\begin{figure}[h]
\centering
\includegraphics[width=\textwidth]{differential_t2_readout.png}
\caption{2T differential GMOS readout schematic. Each GMOS drain is shorted
  to its gate (2T method). A series bias resistor $R_s$ connects each drain
  to the B2961A supply. The drain nodes feed the LIA differential input
  (A and B). The LIA reference sine output drives both heater resistors in
  parallel.}
\label{fig:gmos_2t_schematic}
\end{figure}

\subsubsection{Operating Point}

% TODO: fill in measured / target values.
\begin{table}[h]
\centering
\caption{Nominal operating conditions for the GMOS readout experiment.}
\label{tab:op_point}
\begin{tabular}{lll}
\toprule
Parameter & Symbol & Value \\
\midrule
B2961A supply voltage    & $V_{DD}$               & \ldots\,V \\
Series bias resistor     & $R_s$                  & $\geq 500\,\Omega$ \\
Quiescent drain voltage  & $V_{DS,Q}$             & \ldots\,V \\
Quiescent drain current  & $I_{D,Q}$              & \ldots\,$\mu$A \\
Heater drive amplitude   & $V_{\mathrm{ref}}$     & \ldots\,V$_{\mathrm{pp}}$ \\
Heater drive frequency   & $f_{\mathrm{ref}}$     & \ldots\,Hz \\
LIA time constant        & $\tau$                 & \ldots\,s \\
LIA filter order         & ---                    & --- \\
\bottomrule
\end{tabular}
\end{table}


\end{document}
